\subsection{区间}

设 E 是一个有序集，并设 a, b 是 E 的两个元素，使得 \(a \leqslant b\) 。E 中由满足 \(a \leqslant x \leqslant b\) 的元素 x 组成的子集称为以 a 为左端点、b 为右端点的闭区间，并记为 {[}a, b{]}。所有满足 \(x \in E\) 且 \(a \leqslant x < b\) （相应地 \(a < x \leqslant b\) ）的 x 所成的集合称为以 a 和 b 为端点的右半开（相应地，左半开）区间，并记为 {[}a, b{[} （相应地 {]}a, b{]}）。所有满足 \(x \in E\) 且 a \textless{} x \textless{} b 的 x 所成的集合称为以 a 和 b 为端点的开区间，并记为 {]}a, b{[}。注意，闭区间永远非空；

另一方面，闭区间 \([a, a]\) 是集合 \(\{a\}\) ；而区间 \([a, a[, ]a, a], ]a, a[\) 都是空的；且开区间 \(]a, b[\) 即使当 \(a < b\) 时也可能为空。

设 \(a\) 为 \(E\) 的一个元素。所有满足 \(x \in E\) 且 \(x \leqslant a\) （相应地 \(x < a\) ）的 x 所成的集合称为左无界、右端点为 \(a\) 的闭（相应地，开）区间，并记作 \(] \leftarrow, a]\) （相应地 \(] \leftarrow, a[\) ）；同样地，所有满足 \(x \in E\) 且 \(x \geqslant a\) （相应地 \(x > a\) ）的 x 所成的集合称为左端点为 \(a\) 、右无界的闭（相应地，开）区间，并记为 \([a, \leftarrow[\) （相应地 \(]a, \leftarrow[\) ）。最后， \(E\) 本身可以看作左、右均无界的开区间，记为 \(] \leftarrow, \rightarrow [.\)

命题 13. 在一个格中，两个区间的交集是一个区间。例如考虑两个闭区间 \([a, b]\) 与 \([c, d]\) 的交集，并设 \(\alpha = \sup (a, c)\) , \(\beta = \inf (b, d)\) 。如果 \(a \leqslant x \leqslant b\) 与 \(c \leqslant x \leqslant d\) 两者都成立，那么我们有 \(\alpha \leqslant x \leqslant \beta\) ，且反之亦然；如果 \(\alpha \not\leqslant \beta\) ，则 \([a, b]\) 与 \([c, d]\) 的交集为空；如果 \(\alpha \leqslant \beta\) ，则此交集为 \([\alpha, \beta]\) 。其余情形下的证明留给读者完成。

